\documentclass[fleqn,10pt]{wlscirep}
\usepackage[utf8]{inputenc}
\usepackage[T1]{fontenc}
\usepackage{amsmath,amssymb,graphicx,booktabs,array,longtable}
\usepackage{placeins}
\usepackage{url}
\renewcommand{\topfraction}{.95}
\renewcommand{\textfraction}{.05}
\renewcommand{\floatpagefraction}{.8}
\newcommand{\FixedCov}{2.5920}
\newcommand{\OptCov}{4.3439}
\newcommand{\AbsGain}{1.7519}
\newcommand{\RelGain}{67.59}
\newcommand{\BlindGain}{1.80}
\newcommand{\MCOpt}{4.3515}
\newcommand{\MCLow}{4.3204}
\newcommand{\MCHigh}{4.3827}
\newcommand{\PairGain}{1.7825}
\newcommand{\PairLow}{1.7426}
\newcommand{\PairHigh}{1.8224}
\newcommand{\OceanSCSFixed}{2.2816}
\newcommand{\OceanSCSSynthetic}{3.6065}
\newcommand{\OceanSCSOpt}{3.8180}
\newcommand{\OceanSCSGain}{1.5364}
\newcommand{\OceanSCSGainLow}{1.3004}
\newcommand{\OceanSCSGainHigh}{1.7488}
\newcommand{\OceanKuroFixed}{1.6995}
\newcommand{\OceanKuroSynthetic}{2.4465}
\newcommand{\OceanKuroOpt}{2.3870}
\newcommand{\OceanKuroGain}{0.6875}
\newcommand{\OceanKuroGainLow}{0.5054}
\newcommand{\OceanKuroGainHigh}{0.8794}
\newcommand{\OceanNAFixed}{0.4060}
\newcommand{\OceanNASynthetic}{0.4724}
\newcommand{\OceanNAOpt}{0.7041}
\newcommand{\OceanNAGain}{0.2981}
\newcommand{\OceanNAGainLow}{0.1249}
\newcommand{\OceanNAGainHigh}{0.4852}
\newcommand{\E}{\mathbb{E}}
\newcommand{\Prob}{\mathbb{P}}
\title{Probabilistic Coverage Optimization for Underwater Sensor Networks under Three-Dimensional Localization Uncertainty}
\author[1,*]{Kong Sibo}
\author[2]{Hou Xinyu}
\author[1]{Jin Guanghao}
\author[1]{Zhang Ping}
\author[2]{Zhang Shengti}
\affil[1]{Beijing Polytechnic University, Beijing 100176, China}
\affil[2]{Beijing Information Science and Technology University, Beijing 100192, China}
\affil[*]{Corresponding author: kongsibo@bpu.edu.cn}
\keywords{Underwater sensor networks, Probabilistic coverage, Localization uncertainty, Truncated normal distribution, Powell optimization, Monte Carlo validation}
\begin{abstract}
Underwater sensor deployment is sensitive to three-dimensional positioning uncertainty and physical depth constraints. We formulate coverage by calculating the sensor-union probability at each target point before averaging over the monitoring volume. Physically truncated coordinates and spherical sensing are combined with bounded Powell optimization of a nine-sensor grid. Reproducible experiments include quadrature refinement, independent Monte Carlo validation, matched particle-swarm and artificial-fish-swarm comparisons, and sensitivity analyses. In a declared benchmark with sensing radius 20 m and horizontal and vertical parent scales 120 and 40 m, expected coverage increases from \FixedCov\% to \OptCov\%. Thirty independent Monte Carlo batches estimate optimized coverage as \MCOpt\%, with a 95\% interval of [\MCLow\%, \MCHigh\%]. A further experiment uses real Copernicus Marine GLORYS12V1 currents at three ocean sites, with 2018 for planning and 2019 held out. Shared environmental forcing is handled through conditional union probabilities and explicitly assumed drift and residual-error models. The study quantifies coverage gains, optimizer variability and sensitivity to environmental change, while identifying the resolution, calibration and actuation evidence needed for practical deployment.
\end{abstract}
\begin{document}
\flushbottom\maketitle\thispagestyle{empty}

\section*{Introduction}
Underwater sensor networks (UWSNs) are widely used in marine environmental monitoring, ecological observation, subsea infrastructure inspection and offshore resource observation\cite{akyildiz2005}. Unlike terrestrial wireless sensor networks, UWSNs operate in complex three-dimensional environments where sensor positions are difficult to determine precisely. Ocean currents, buoyancy variation, acoustic communication delay and localization errors can cause deployed nodes to deviate from their intended positions. As a result, each node's effective sensing region is uncertain, and deterministic geometry alone may inadequately describe monitoring performance.

Sensor deployment is therefore a key issue in UWSN design. Deterministic coverage analysis establishes whether a region is covered by sensors with known locations\cite{huang2003}, while deployment and target-localization formulations connect sensing geometry to placement decisions\cite{zou2004}. These are useful reference models, but their assumptions must be revisited when location errors are comparable to a sensor's sensing radius. Uncertainty-aware deployment has consequently been investigated explicitly\cite{zouuncertainty2004}. In underwater settings, stochastic model refinement is also relevant to positioning itself\cite{zhao2020}. A probabilistic treatment does not replace the physical detection event: it averages that same event over uncertainty in actual sensor positions.

This study focuses on probabilistic coverage optimization under three-dimensional localization uncertainty. The horizontal coordinates have Gaussian parent distributions, and absolute depth has a truncated normal distribution constrained by the feasible water column. The model distinguishes the nominal deployment, the realized position and their difference. It also distinguishes the monitoring region from the larger admissible positioning region. These distinctions matter because a sensor outside the target volume can still cover part of it, whereas a nominal point inside a physical boundary does not ensure that a Gaussian realization respects that boundary.

The nonlinear objective is evaluated by numerical integration. A derivative-free Powell method\cite{powell1964} searches the deployment parameters, and independent Monte Carlo trials\cite{metropolis1949} validate the same expected covered-volume fraction. The central requirement is consistency between the pointwise sensing event, the network union event and the quantity averaged over the monitoring region. A product of already spatially averaged probabilities generally fails to preserve that relationship.

The study addresses three research questions. \textbf{RQ1:} How can three-dimensional localization uncertainty be incorporated into a physically constrained probabilistic coverage model? \textbf{RQ2:} How do horizontal and vertical uncertainties affect coverage under different deployment configurations? \textbf{RQ3:} Can derivative-free optimization improve coverage, and how does its computational performance compare with alternative methods under matched conditions?

The contributions are a consistent pointwise-to-volume formulation, a reproducible spherical-probability evaluator, an explicitly constrained grid-deployment problem, and a controlled numerical study. The latter includes fixed-layout and single-variable ablations, 20 runs per optimizer, 30 independent validation batches, spatial coverage maps and sensitivity to uncertainty and sensing radius. A real-current replay further uses three public Copernicus Marine sites with a chronological planning/test split and common environmental forcing. The framework retains an explicit route towards movement-aware deployment, while separating demonstrated results from capabilities requiring additional implementation.

\section*{Related Studies}
\subsection*{Coverage Modeling under Localization Uncertainty}
Deterministic coverage models provide a geometric starting point\cite{huang2003,zou2004}. Once sensor positions become random, coverage at a target becomes a probability, and location uncertainty must be integrated over the sensing event\cite{zouuncertainty2004}. A key distinction is whether uncertainty belongs to the target, to the sensor position, or to detection conditional on distance. Here the target is a spatial integration variable, sensor position is random, and detection is the binary event that their Euclidean distance does not exceed a specified radius. This interpretation makes the measured quantity unambiguous.

The Gaussian parent distribution supplies a parsimonious model of residual positioning error. It is a declared modeling assumption, motivated by the aggregation of multiple perturbations, rather than a universal empirical law for underwater localization. Physical truncation excludes impossible actual positions. Stochastic positioning studies emphasize the importance of refining uncertainty descriptions\cite{zhao2020}; anisotropy, bias, correlation and time variation can all be important. The independent, homogeneous benchmark used here provides a controlled starting point whose limitations are subsequently examined.

\subsection*{Optimization Methods for Sensor Deployment}
Powell's method constructs search directions and performs derivative-free line searches\cite{powell1964}. It is appropriate for a low-dimensional objective that is deterministic once the quadrature rule is fixed. It does not establish global optimality and may depend on initialization. Particle swarm optimization (PSO) instead maintains a population influenced by individual and collective best positions\cite{kennedy1995,clerc2002}. Artificial fish swarm algorithms (AFSA) use behaviors such as following, swarming and local searching\cite{li2002}. Their performance depends on implementation, population size, parameters and evaluation budget. Accordingly, the present comparison reports explicit variants and repeated runs, without attributing unsupported universal strengths or weaknesses to these algorithm families.

Learning-based optimization and underwater localization are complementary research directions. A reinforcement-learning multi-swarm optimizer addresses large-scale optimization\cite{wang2024}, while underwater acoustic localization and robotic tracking have also been studied using reinforcement learning\cite{zhou2024,masmitja2023}. These studies motivate future adaptive systems; they do not provide controlled coverage baselines for the present grid-deployment objective. In particular, a tracking paper is not used as a citation for PSO.

\subsection*{Probabilistic Validation and Research Gap}
Monte Carlo integration estimates expectations using independent random observations\cite{metropolis1949,papoulis2002}. To validate expected coverage, each observation must combine a target uniformly sampled from the same monitoring volume and an independently realized sensor network. Sampling only sensor positions at a selected point would validate a different quantity. Spatial grid values are likewise deterministic evaluations, not independent experimental observations. The validation design therefore distinguishes optimizer variability, Monte Carlo uncertainty and quadrature refinement.

\subsection*{Summary of Gaps and Positioning of This Study}
The present study connects a physically constrained three-dimensional location model, pointwise union coverage, bounded optimization and independent same-target validation within one reproducible workflow (Fig.~\ref{fig:framework}). It does not claim that Gaussian uncertainty, Powell optimization or Monte Carlo simulation is new in isolation. Its purpose is to establish consistent calculation and interpretable deployment comparisons, including situations in which an optimized layout loses its advantage when uncertainty changes.

\begin{figure}[!t]\centering
\includegraphics[width=\linewidth]{figures/fig1_framework.pdf}
\caption{\textbf{Probabilistic deployment framework.} The physical model, pointwise union probability, feasible grid and numerical optimizer feed an independent validation and analysis stage. All implemented benchmark paths use the same positive-downward depth convention and physical bounds. The extension interface identifies how time-varying environmental models could supply distributions to the coverage evaluator; it does not denote an implemented CFD--LSTM controller.}\label{fig:framework}
\end{figure}

\section*{Materials and Methods}
\subsection*{Problem Formulation}
Let $\Omega\subset\mathbb R^3$ be the monitoring volume, $\mathbf q=(q_x,q_y,q_z)$ a target point and $\mathbf s_i$ the nominal position of sensor $i$. The actual position is the random vector $\mathbf S_i$. Absolute depth is positive downward throughout. The additive error is defined only after the actual-position model is specified, as $\boldsymbol\varepsilon_i=\mathbf S_i-\mathbf s_i$. Thus a nominal depth of 150 m is not a 150 m mean additive error.

The declared computational benchmark uses
\begin{equation}\label{eq:domains}
\Omega=[-50,50]\times[-10,10]\times[137.5,162.5]\ \mathrm{m}^3,
\qquad \mathcal F=[-500,500]^2\times[120,270]\ \mathrm{m}^3.
\end{equation}
The dimensions of $\Omega$ are $L=100$ m, $B=20$ m and $H=25$ m, giving $|\Omega|=50{,}000$ m$^3$. The larger region $\mathcal F$ bounds nominal and actual sensor positions. The sensing radius is $R=20$ m. These values define a synthetic benchmark, not a measurement campaign or a calibrated ocean site. Table~\ref{tab:parameters} specifies the remaining parameters.

For coordinate $k\in\{x,y,z\}$, let $(a_k,b_k)$ be the corresponding bounds in $\mathcal F$ and let $\sigma_k>0$ be the parent Gaussian scale. The marginal density is
\begin{equation}\label{eq:density}
 f_{ik}(u)=\frac{\phi((u-s_{ik})/\sigma_k)}{\sigma_k Z_{ik}}\,\mathbf 1_{[a_k,b_k]}(u),
 \qquad Z_{ik}=\Phi((b_k-s_{ik})/\sigma_k)-\Phi((a_k-s_{ik})/\sigma_k),
\end{equation}
where $\phi$ and $\Phi$ denote the standard normal density and cumulative distribution function. We set $\sigma_x=\sigma_y=\sigma_h=120$ m and $\sigma_z=40$ m. These are standard deviations of the parent distributions, not variances and not generally the standard deviations after truncation. Horizontal truncation is included so that physical bounds apply to every actual coordinate. For central layouts its effect is small, but it is included in both integration and sampling.

Depth truncation changes the mean. In particular,
\begin{equation}\label{eq:mean}
\E[S_{iz}]=s_{iz}+\sigma_z\frac{\phi((a_z-s_{iz})/\sigma_z)-\phi((b_z-s_{iz})/\sigma_z)}{Z_{iz}}.
\end{equation}
Consequently $\E[\varepsilon_{iz}]$ need not vanish. This is crucial near the shallow bound: placing the Gaussian parent centre at 120 m yields an actual mean depth near 152 m under the benchmark scale. The optimization acts on the nominal parent-centre parameter. A controller whose commanded depth represents the mean actual depth would require a different parameterization.

In the synthetic benchmark, the coordinates of an individual sensor are assumed independent, and different sensors have independent realizations. The environmental replay below instead assumes independence conditional on a shared current state. Their joint density therefore factorizes as $f_i(\mathbf u)=\prod_k f_{ik}(u_k)$\cite{papoulis2002}. These are modeling assumptions, not conclusions obtained from the supplied data. Shared current-induced displacement or localization-reference errors can violate sensor independence. In that situation the network probability must be evaluated from a joint distribution rather than from a product of individual complements.

\begin{table}[htbp]\centering\small
\caption{\textbf{Complete benchmark and numerical settings.} The same physical configuration and objective are used by all optimizers. These are the synthetic-benchmark settings. The real-current replay retains the physical geometry and residual scales and adds the environmental inputs described below. Parent Gaussian scales are stated in metres.}\label{tab:parameters}
\begin{tabular}{p{.29\linewidth}p{.64\linewidth}}\toprule
Item & Setting\\\midrule
Target / physical region & $\Omega=[-50,50]\times[-10,10]\times[137.5,162.5]$ m$^3$; $\mathcal F=[-500,500]^2\times[120,270]$ m$^3$\\
Sensors and sensing & $N=9$; isotropic binary sphere; $R=20$ m\\
Position distribution & Independent truncated Gaussian coordinates and independent sensors; $\sigma_h=120$ m, $\sigma_z=40$ m\\
Nominal geometry & Planar $3\times3$ grid; centre $(c_x,c_y)$, pitch $d$, common depth $z_0$\\
Optimization bounds & $40\le d\le500$ m, $120\le z_0\le270$ m, $|c_x|,|c_y|\le500-d$ m\\
Fixed baseline & $(c_x,c_y,d,z_0)=(0,0,100,150)$ m; the coordinate pitch is 100 m\\
Nominal separation & Minimum 40 m; no constraint is imposed between realized random positions\\
Training quadrature & 16-point depth integration; $5\times3\times5$ target Gauss--Legendre nodes\\
Refinement levels & $(24;7,5,7)$, $(48;11,7,15)$, $(80;15,9,25)$; first entry is depth order\\
Powell settings & Unit-box variables; coordinate directions; $\mathrm{xtol}=10^{-5}$, $\mathrm{ftol}=10^{-9}$; 100 iterations; at most 600 evaluations\\
Repeated optimization & 20 seeds, 2100--2119; PSO and AFSA population 20; all calls counted\\
Monte Carlo validation & 30 seeds, 31000--31029; 50,000 independent target--array trials per seed; paired layouts\\
Sensitivity / design sample & $\sigma_h\in\{30,60,120,180\}$ m, $\sigma_z\in\{0,15,40,60\}$ m; 600 independent uniform designs, seed 43210\\\bottomrule
\end{tabular}
\end{table}

\subsection*{Probabilistic Coverage Modeling}
A sensor covers $\mathbf q$ when $\|\mathbf S_i-\mathbf q\|_2\le R$. The corresponding probability is
\begin{equation}\label{eq:single}
p_i(\mathbf q)=\int_{\mathcal F}\mathbf 1_{\{\|\mathbf u-\mathbf q\|_2\le R\}}f_i(\mathbf u)\,d\mathbf u.
\end{equation}
This integrates over a sphere intersected with the feasible physical region. Integrating over an axis-aligned rectangle enlarged by $R$ would count points outside the sensing sphere and would represent a different event.

Under independent sensor positions, the network union probability and expected covered-volume fraction are
\begin{align}
P_{\mathrm{cov}}(\mathbf q;\mathcal S)&=1-\prod_{i=1}^{N}\{1-p_i(\mathbf q)\},\label{eq:union}\\
C(\mathcal S)&=\frac{1}{|\Omega|}\int_{\Omega}P_{\mathrm{cov}}(\mathbf q;\mathcal S)\,d\mathbf q.\label{eq:coverage}
\end{align}
Here $\mathcal S=\{\mathbf s_i\}_{i=1}^N$. Fubini's theorem also gives $C=\E[|\Omega\cap\bigcup_i B(\mathbf S_i,R)|]/|\Omega|$. Thus $C$ is an expected fraction of volume, not the probability that the entire monitoring region is covered. Its complement $1-C$ is the expected uncovered fraction, not a thresholded ``blind-area'' percentage unless a separate threshold is defined.

The order of union formation and spatial averaging cannot be interchanged. For two sensors, with $\bar p_i=|\Omega|^{-1}\int_\Omega p_i(\mathbf q)d\mathbf q$,
\begin{equation}\label{eq:covariance}
C=\bar p_1+\bar p_2-\E_{\mathbf Q}[p_1(\mathbf Q)p_2(\mathbf Q)],\qquad
\{1-(1-\bar p_1)(1-\bar p_2)\}-C=\operatorname{Cov}_{\mathbf Q}(p_1,p_2).
\end{equation}
Even when sensor errors are independent, their probability fields need not be uncorrelated over a common target. This explains why independence across sensors does not justify multiplying spatially averaged probabilities. Supplementary Table S4 quantifies the difference in the current configurations and also reports the sum of single-sensor means, which is only a union bound.

\subsection*{Numerical Integration and the Zero-Depth-Error Limit}
The benchmark has equal horizontal parent scales. Before horizontal truncation, squared horizontal distance divided by $\sigma_h^2$ follows a noncentral chi-square distribution with two degrees of freedom and noncentrality
\begin{equation}\label{eq:noncentrality}
\lambda_i(\mathbf q)=\frac{(q_x-s_{ix})^2+(q_y-s_{iy})^2}{\sigma_h^2}.
\end{equation}
All sensing spheres around target points in the reported radius range lie inside the horizontal physical bounds. Consequently horizontal conditioning contributes only its normalization factor, and equation~\eqref{eq:single} reduces to the one-dimensional integral
\begin{equation}\label{eq:quadrature}
p_i(\mathbf q)=\frac{1}{Z_{ix}Z_{iy}}\int_{\ell(\mathbf q)}^{u(\mathbf q)}
F_{\chi'^2_2(\lambda_i)}\!\left(\frac{R^2-(z-q_z)^2}{\sigma_h^2}\right)f_{iz}(z)\,dz,
\end{equation}
where $\ell=\max(a_z,q_z-R)$ and $u=\min(b_z,q_z+R)$. The probability is zero if $u\le\ell$. This reduction is specific to isotropic horizontal uncertainty and the stated geometric condition. More general anisotropy or intersection with a horizontal wall requires direct integration over the clipped horizontal disk.

We apply Gauss--Legendre quadrature to the depth integral and tensor-product Gauss--Legendre quadrature to the target volume. Normalized target weights sum to one. At each target node the sensor probabilities are combined using equation~\eqref{eq:union}, then weighted and summed. Refinement is reported separately from Monte Carlo confidence intervals. The comparison between finite rules is a convergence diagnostic, not a formal mathematical error bound.

For $\sigma_z=0$, actual depth equals nominal depth and the vertical distribution is degenerate. We evaluate the horizontal disk probability directly, with squared disk radius $\max\{R^2-(q_z-z_0)^2,0\}$. We retain the same three-dimensional target volume in this ablation, so only vertical uncertainty changes. A two-dimensional target-plane experiment would have a different denominator and is not directly interchangeable with the reported volume coverage. The zero-depth-error analysis therefore uses a consistent target definition.

\subsection*{Optimization Framework}
Nominal nodes follow the prescribed grid
\begin{equation}\label{eq:grid}
\mathbf s_{ij}=(c_x+i d,c_y+j d,z_0),\qquad i,j\in\{-1,0,1\}.
\end{equation}
The minimum nominal separation is $d\ge40$ m. For bounded search variables $\boldsymbol\theta=(\theta_1,\ldots,\theta_4)\in[0,1]^4$, we use
\begin{align}
d&=40+460\theta_3,&z_0&=120+150\theta_4,\nonumber\\
c_x&=(2\theta_1-1)(500-d),&c_y&=(2\theta_2-1)(500-d).\label{eq:parameterization}
\end{align}
Every nominal node is consequently feasible without a penalty function. The problem is $\max_{\boldsymbol\theta\in[0,1]^4}C(\mathcal S(\boldsymbol\theta))$. This optimizes four grid parameters, not 27 independent coordinates. Nor does nominal separation guarantee collision avoidance of realized uncertain positions. The present static sensing model contains no contact mechanics, communication-connectivity constraint or depth-dependent acoustic propagation law.

The bounded modified Powell implementation uses a deterministic training quadrature, identity initial directions and the stopping settings in Table~\ref{tab:parameters}. Twenty independent initial points test sensitivity to starting location. The best Powell layout is selected by training coverage and then evaluated on finer rules and independent Monte Carlo trials. No Monte Carlo observations are used to select a layout. Exact initial points, evaluation histories, stop messages, coordinates and software versions are saved with the experiment data.

\subsection*{Matched Baselines and Algorithm Budgets}
The fixed nine-node grid has $d=100$ m and $z_0=150$ m. A pitch-only layout sets $d=40$ m while retaining $z_0=150$ m; a depth-only layout sets $z_0=120$ m while retaining $d=100$ m. These ablations separate the geometric and depth-model contributions to improvement. A single sensor at $(0,0,150)$ m remains an interpretive reference for the benefit of additional sensors. It is not used as evidence that an optimizer is superior.

Powell, PSO, the specified AFSA variant and uniform random search receive the same maximum budget of 600 objective calls, the same domain, sensing radius, node count, uncertainty law, feasible parameterization and training rule. For each seed, Powell starts at the first member of the same randomly generated population used to initialize PSO and AFSA. Powell may terminate before exhausting its allowance; actual counts are reported. Population methods include their initialization evaluations in the budget, and AFSA counts exploratory and centroid evaluations, including candidates that are not accepted.

PSO uses 20 particles, inertia 0.7298, cognitive and social coefficients 1.49618, and a velocity cap of 0.2 in each unit-box coordinate. Positions are clipped to the box and outward velocities are set to zero. The AFSA variant uses 20 fish, Euclidean neighborhoods in the unit box, following and centroid candidates subject to a crowding fraction of 0.6, up to three prey trials, and random movement when no improved candidate is found. Visual radius decreases linearly from 0.55 to 0.05, and movement step from 0.11 to 0.01 as the evaluation budget is consumed. A best-observed bulletin retains every evaluated candidate, including exploration points. The full update rules are given in Supplementary Methods S1 and the executable source; conclusions apply to this variant and parameterization.

Each algorithm is run with seeds 2100--2119. Final candidates are reevaluated with the fine rule, and comparisons report the distribution of those final values. Training convergence curves show the median and interquartile range of the best value seen by each call count. For an early-stopped Powell run its final value is carried forward only for plotting. Measured runtime covers optimization calls and algorithm overhead, excluding the final refined evaluation; it is specific to the recorded computer and software, not an operational deployment time.

\subsection*{Independent Monte Carlo Validation and Statistical Units}
For repetition $r$, each trial $m$ draws $\mathbf Q_{rm}$ uniformly from $\Omega$, then draws the full sensor array independently from the model in equation~\eqref{eq:density}. The Bernoulli observation is
\begin{equation}\label{eq:mc}
I_{rm}=\mathbf 1\{\min_i\|\mathbf S_{irm}-\mathbf Q_{rm}\|_2\le R\},\qquad
\widehat C_r=\frac{1}{M}\sum_{m=1}^M I_{rm}.
\end{equation}
Since $\E[I_{rm}]=C$, this validates equation~\eqref{eq:coverage} directly. We use $M=50{,}000$ trials and 30 independent repetitions, totaling 1.5 million trials per layout. Independent uniform variates transformed through each truncated normal inverse CDF generate actual positions. Random seeds and bit-packed trial outcomes are archived.

The same targets and underlying independent uniform variates are reused across layouts within a repetition. This common-random-number design pairs configurations without creating dependence between sensors in one configuration or between trials. For each layout, the reported standard deviation is the sample standard deviation of the 30 batch estimates. An approximate 95\% interval is $\bar C\pm t_{0.975,29}s_C/\sqrt{30}$. The paired gain interval is calculated from the 30 within-repetition differences between optimized and fixed coverage, not from an unpaired combination of intervals.

Optimizer runs and Monte Carlo repetitions are separate statistical units. Coordinate sweeps, spatial maps and quadrature nodes are deterministic evaluations. A further 600 uniformly sampled feasible grid centres and nominal depths, with fixed pitch 40 m, describe how performance varies across a design region; their boxplots show conditional spread across designs, not uncertainty around an optimum. No hypothesis-test significance, Pearson significance or regression goodness-of-fit is inferred from deterministic mesh values.


\subsection*{Real Ocean Data and Current-Conditioned Replay}
To examine environmental relevance, we use real daily currents from the Copernicus Marine Global Ocean Physics Reanalysis (GLORYS12V1; product DOI 10.48670/moi-00021)\cite{cmems_glorys,lellouche2021}. The original manuscript named CMEMS without specifying a product, site, date or raw file. We now identify the complete source explicitly: Mercator Ocean International's public OceanBench historical GLORYS12V1 derivative, distributed on a $1^\circ$ horizontal grid\cite{oceanbench_archive}. This is observation-constrained ocean reanalysis, not a set of direct sensor-position or current-meter measurements. The underlying product has a native $1/12^\circ$ grid and 50 depth levels, but the downloaded derivative has $1^\circ$ resolution; the native resolution is not assigned to these experiments.

We select South China Sea $(115.5^\circ\mathrm E,20.5^\circ\mathrm N)$, Kuroshio vicinity $(130.5^\circ\mathrm E,25.5^\circ\mathrm N)$ and North Atlantic $(44.5^\circ\mathrm W,40.5^\circ\mathrm N)$ cells before optimization. Eastward and northward velocities at the retained depths 130.6660 and 155.8507 m are linearly interpolated to 150 m. All 730 daily records per site in 2018--2019 are finite and unique; no missing dates are filled and no temporal interpolation is applied. Original metadata, exact public requests, source-chunk hashes, retained-depth NetCDF subsets and the extracted 2,190-row CSV are archived. The upstream release identifier is not specified in the public archive; its exact files are identified by request and checksum rather than assigned the version of the current catalogue.

The 2018 record is used for planning, and the complete 2019 record is held out. With a prespecified horizon $\Delta t=600$ s, a frozen-current displacement proxy is
\begin{equation}\label{eq:current_drift}
\mathbf b_t=\Delta t\,(u_t,v_t),\qquad
\boldsymbol\mu_{it}=\mathbf s_i+(b_{tx},b_{ty},0).
\end{equation}
Here $u_t,v_t$ are real reanalysis velocities in m s$^{-1}$. The conversion to displacement, uniform forcing across the small target region, and residual parent scales $(120,120,40)$ m are modeling assumptions. They are not fitted sensor-localization errors or observed trajectories. Conditional actual coordinates use equation~\eqref{eq:density} with parent centre $\boldsymbol\mu_{it}$ and the same physical bounds. Nominal nodes remain feasible, whereas a drifted parent centre can lie outside the box; the conditioned actual distribution remains physically bounded. This static distributional replay does not model wall-contact dynamics or a continuous trajectory.

Every sensor shares the current day. Residual errors are independent conditional on that day, so the environmental objective is
\begin{equation}\label{eq:current_objective}
C_{\rm env}(\mathcal S)=\frac1T\sum_{t=1}^{T}\frac1{|\Omega|}
\int_\Omega\left[1-\prod_{i=1}^{N}\{1-p_i(\mathbf q\mid\mathbf b_t)\}\right]d\mathbf q.
\end{equation}
The union is formed before both spatial and current-state averaging. The common forcing induces unconditional sensor dependence; replacing the conditional product by a product of current-averaged sensor probabilities would lose that dependence. Equation~\eqref{eq:quadrature} evaluates each conditional sensing probability with its displaced horizontal centre and recomputed physical normalizers.

For economical planning, 16 representative current states are obtained by k-means++ using only the 365 training displacements, seed 6200 and 60 centroid iterations. Final nearest-centre assignments give empirical weights,
\begin{equation}\label{eq:current_reduction}
w_k=\frac{n_k}{T_{\rm train}},\qquad
C_K(\mathcal S)=\sum_{k=1}^{K}w_k C(\mathcal S\mid\bar{\mathbf b}_k),\qquad K=16.
\end{equation}
All four optimizers use this same weighted objective, nine nodes, feasible grid and maximum 300 calls, with 20 seeds 5100--5119 per site. The planning rule uses 12 depth nodes and $3\times2\times3$ target nodes. Final candidates are reevaluated on the 16 training states using 32 depth nodes and $9\times5\times11$ target nodes. The highest refined training score selects one deployment per method; no 2019 value is used for selection. Full-year training evaluations and 32-state evaluations quantify state-reduction error separately from quadrature refinement.

All selected deployments are evaluated on all 365 held-out current days with 48 depth nodes and $11\times7\times15$ target nodes. Powell test results are additionally refined to $(64;13,9,19)$. Fixed, synthetic-optimum and current-aware Powell layouts receive 30 independent Monte Carlo batches of 50,000 trials at each site. Each trial draws one held-out day uniformly, one uniform target and independent residual variates for all nine nodes; the same day is shared across nodes. Common random numbers pair layouts within a batch. The 13.5 million trial outcomes validate equation~\eqref{eq:current_objective} conditional on the empirical record.

Daily values form an autocorrelated environmental time series. We therefore use a circular moving-block bootstrap for temporal variation of the paired gain $g_t$, with 14-day blocks, 2,000 draws and seed 7100:
\begin{equation}\label{eq:temporal_bootstrap}
\widehat G_b^*=\frac1T\sum_{j=1}^{T}g_{I_{bj}},\qquad
\mathrm{CI}_{\rm temporal}=[Q_{0.025}(\widehat G^*),Q_{0.975}(\widehat G^*)].
\end{equation}
Blocks preserve local serial dependence; 7- and 30-day block lengths are also archived as a sensitivity check. These intervals describe variation in the finite environmental record under block resampling. Monte Carlo intervals describe independent target--array sampling conditional on that record, and deterministic quadrature changes describe numerical refinement. None establishes climate-independent performance or a sea-trial confidence interval.

\subsection*{Formula-to-Code Reproducibility}
The separate Formula, Pseudocode and Source Code Companion covers all displayed equations and computational inline expressions through stable identifiers F01--F27. Each entry gives units, ordered pseudocode and exact executable file/function locations generated from the Python syntax tree. The reference functions and production functions are linked in \texttt{formula\_registry.json}. Independent checks verify density normalization, actual means, the spherical integral, aggregation order, physical sampling support, zero-current equivalence, percent units and shared-forcing semantics. Proposed energy interfaces are explicitly identified and reject missing calibrated inputs.

The accompanying figure/table provenance inventory identifies every input and plotted quantity. All environmental experiment inputs are downloaded public reanalysis records. Numerical-benchmark samples are explicitly mathematical checks rather than environmental observations, and Monte Carlo outcomes are computed integration variables rather than measured ocean data. No assumed position-error scale is represented as estimated from the current dataset.

\section*{Results}
\subsection*{Solving the Coverage Model under Three-Dimensional Localization Error}
\subsubsection*{Maximum Coverage Analysis under Zero Depth Error}
The error distributions and the three-dimensional coverage field are shown in Fig.~\ref{fig:distribution}. A density has units of inverse length, while a coverage probability is dimensionless; the two are distinguished in axis labels and captions. Horizontal uncertainty spreads possible sensor locations far beyond the small monitoring volume in this benchmark. Its scale of 120 m is six times the sensing radius, so low absolute coverage is expected even for nine sensors.

With zero vertical uncertainty, the fixed nine-node grid gives approximately 6.6815\% coverage at $\sigma_h=120$ m. Introducing $\sigma_z=40$ m at that same nominal geometry reduces coverage to \FixedCov\%. The comparison keeps the target volume and all other parameters fixed and therefore isolates the effect of vertical uncertainty. The magnitude cannot be interpreted as a universal penalty for ``3D modeling'': it depends on the nominal depth, physical truncation and target geometry.

The zero-depth-error case also clarifies why an optimized location must be interpreted together with its uncertainty law. A nominal depth near 120 m performs well with the broad truncated depth distribution, but with zero vertical uncertainty most of the target volume is outside its vertical sensing range. The corresponding optimized-for-40-m layout covers only approximately 0.1358\% when $\sigma_z=0$. Thus there is no general guarantee that decreasing a scale parameter improves the performance of a layout selected under a different distribution.

\begin{figure}[htbp]\centering
\includegraphics[width=\linewidth]{figures/fig2_distributions_field.pdf}
\caption{\textbf{Position distributions, coverage field and sampled deployment space.} (a) Horizontal error density with a 10,000-draw synthetic histogram (seed 7654). (b) Actual-depth densities for nominal depths of 150 and 120 m, both using parent scale 40 m and physical bounds [120,270] m. (c) Optimized pointwise coverage at target depth 150 m. This is a probability surface, not a temperature field. (d) Coverage of 600 uniformly sampled grid centres and depths at pitch 40 m; colors encode the computed objective. The spread is descriptive and does not provide a confidence interval for the optimized configuration.}\label{fig:distribution}
\end{figure}

\subsubsection*{Sensitivity Analysis of Coverage under Positioning Errors}
Figure~\ref{fig:sensitivity} retains the horizontal, vertical and observation-range analyses while making the varied quantity explicit. For the fixed layout at $\sigma_z=40$ m, increasing $\sigma_h$ from 30 to 60, 120 and 180 m changes coverage from approximately 5.2560\% to 4.0197\%, 2.5920\% and 1.5073\%, respectively. The corresponding values for the selected optimized layout are approximately 22.7498\%, 12.6281\%, 4.3439\% and 2.0932\%. These are sensitivity evaluations at fixed nominal coordinates, not separately reoptimized deployments for each scenario.

Vertical sensitivity is more strongly influenced by truncation. For the fixed layout with $\sigma_h=120$ m, coverage is approximately 4.4309\%, 2.5920\% and 2.0442\% at vertical scales 15, 40 and 60 m. The selected optimized layout gives approximately 3.5704\%, 4.3439\% and 3.4586\%, respectively. The ordering of layouts reverses at 15 m. The depth distribution centered at the shallow boundary places much less probability near the target as its scale shrinks; at intermediate scales it aligns more closely with the target. This explains the non-monotonic behavior without invoking an unmodeled dynamic ocean process.

Increasing the sensing radius increases coverage for any fixed set of actual-position distributions because the sensing spheres are nested. The radius sweep from 10 to 40 m in Fig.~\ref{fig:sensitivity}c demonstrates this property, with all other settings fixed. Depth in Fig.~\ref{fig:sensitivity}d is the nominal absolute depth parameter, not observation distance, sensor-to-target distance or an additive-error mean. Correctly labeling these variables is essential for interpreting the vertical results.

\begin{figure}[htbp]\centering
\includegraphics[width=\linewidth]{figures/fig3_sensitivity.pdf}
\caption{\textbf{Uncertainty and sensing sensitivity.} (a,b) Expected covered-volume percentages for the fixed and selected optimized layouts under 16 combinations of horizontal and vertical parent scales. The zero-vertical-scale row uses a degenerate depth distribution. (c) Radius sweep with the benchmark uncertainty held fixed. (d) Common nominal-depth sweep at the fixed and optimized grid pitches. The optimized coordinates are not reselected for each uncertainty scenario; the panels therefore expose distributional sensitivity rather than a robustness guarantee.}\label{fig:sensitivity}
\end{figure}

\subsubsection*{Multi-Constraint Validation for Sensor Deployment Optimization}
The best selected Powell layout has centre numerically indistinguishable from $(0,0)$ m, pitch approximately 40.00004 m and common nominal depth approximately 120.00000 m. All nominal nodes satisfy the physical box and minimum-separation constraint. Their coordinates and individual spatially averaged contributions are given in Supplementary Table S3; both the fixed and optimized grid are visualized in Supplementary Fig. S3.

The optimum found lies at the lower pitch and depth bounds. It is therefore a boundary solution within the prescribed design family. It does not establish a globally optimal free-form three-dimensional array. The shallow nominal depth is explained by equation~\eqref{eq:mean}; it is not evidence that sensors should always be placed above the monitoring volume. The separation constraint prevents arbitrarily close nominal placement, but actual sensing overlap remains possible and is accounted for through the union formula.

Figure~\ref{fig:constraints} combines coordinate offsets, pitch and spatial maps. Symmetry of the target and the parent distributions explains the matching behavior under positive and negative horizontal offsets. The long target dimension and narrow cross-track dimension nevertheless produce different offset sensitivities. The optimized map remains spatially smooth but does not imply almost complete monitoring. The scales are shared across map panels so that visual differences correspond to probability differences.

\begin{figure}[htbp]\centering
\includegraphics[width=\linewidth]{figures/fig4_constraints_maps.pdf}
\caption{\textbf{Coordinate sensitivity, separation and coverage maps.} (a,b) Horizontal centre offsets applied to the fixed and optimized geometries. (c) Grid-pitch sweep at nominal depths 150 and 120 m. (d--f) Pointwise coverage maps at target depth 150 m for the fixed, pitch-only and fully optimized configurations, shown with a common color scale. These maps preserve spatial information that is lost in a single volume-average metric. Supplementary Fig. S1 provides the corresponding descriptive boxplots across independently sampled designs.}\label{fig:constraints}
\end{figure}

\subsection*{Performance Validation under Multidimensional Constraints}
\subsubsection*{RQ1: Consistency between Probability Modeling and Simulation}
The numerical model and Monte Carlo experiment evaluate the same target--array event. Table~\ref{tab:coverage} reports all five configurations and Supplementary Fig.~S4 compares their estimates. The refined quadrature value falls inside the Monte Carlo interval for each configuration. The optimized estimate is \MCOpt\%, with a between-batch standard deviation of approximately 0.0834 percentage points and interval [\MCLow\%, \MCHigh\%]. This is evidence of consistency under the stated synthetic model, not validation against field observations.

Quadrature refinement changes the numerical answer by much less than the Monte Carlo interval width. For the optimized layout the fine and finest estimates differ by about $3.8\times10^{-8}$ in coverage fraction. Intermediate changes are not necessarily monotonic because the clipped depth interval varies with target depth. The complete refinement results appear in Supplementary Table S2. Repeated Monte Carlo observations quantify sampling error, while refinement compares deterministic integration schemes; these serve different purposes.

\begin{table}[htbp]\centering\small
\caption{\textbf{Controlled layout comparisons and independent validation.} Quadrature uses the finest rule. Monte Carlo (MC) summaries use 30 independent batches of 50,000 trials each. SD is the between-batch standard deviation in percentage points (pp); the interval is for the mean coverage. The single sensor is a node-count reference, while optimization gains use the matched nine-node fixed baseline.}\label{tab:coverage}
\begin{tabular}{lrrrrl}
\toprule
Layout & $N$ & Quadrature (\%) & MC mean (\%) & SD (pp) & 95\% CI (\%) \\
\midrule
Single & 1 & 0.4437 & 0.4443 & 0.0251 & [0.4349, 0.4537] \\
Fixed & 9 & 2.5920 & 2.5690 & 0.0661 & [2.5443, 2.5937] \\
Pitch-only & 9 & 3.6594 & 3.6761 & 0.0865 & [3.6438, 3.7084] \\
Depth-only & 9 & 3.0800 & 3.0668 & 0.0801 & [3.0369, 3.0967] \\
Optimized & 9 & 4.3439 & 4.3515 & 0.0834 & [4.3204, 4.3827] \\
\bottomrule
\end{tabular}

\end{table}



\subsubsection*{RQ2: Coverage Attenuation and Model Dependence}
The sensitivity experiment answers RQ2 at fixed layouts rather than by pooling unrelated configurations. At the fixed baseline, increasing vertical uncertainty reduces probability mass close enough to targets to satisfy the spherical event. At the optimized boundary layout, the relationship is non-monotonic because truncation shifts the actual-depth distribution. Therefore statements about attenuation must name the layout and the parameter range. Supplementary Table S5 preserves the numerical zero-depth and three-dimensional comparison, and the archived CSV files contain the full sensitivity grid.

Spatially averaged coverage also does not determine the probability of a large contiguous blind region. Although the maps characterize pointwise variation, the probability distribution of uncovered connected components would require a further random-field analysis. We report the expected uncovered-volume fraction using $1-C$ rather than infer a thresholded area reduction from colors in a plot.

\subsubsection*{RQ3: Cooperative Gain and the Benefit of Optimization}
The fixed nine-sensor configuration covers \FixedCov\% of the target volume in expectation, compared with approximately 0.4437\% for a single sensor. This node-count comparison illustrates the benefit of additional sensing opportunities but includes resource growth. The direct optimization comparison uses nine sensors on both sides: coverage increases to \OptCov\%, an absolute gain of \AbsGain\ percentage points and relative increase of \RelGain\%.

The paired Monte Carlo gain is \PairGain\ percentage points, with a 95\% interval of [\PairLow,\PairHigh] percentage points. The distinction between relative percent and percentage points avoids overstating the effect. Expected uncovered volume falls from approximately 97.4080\% to 95.6561\%, corresponding to a relative reduction of about \BlindGain\%. This reduction is meaningful but small in absolute monitoring terms; the selected benchmark remains difficult because uncertainty is large relative to the sensing radius.

The pitch-only and depth-only configurations reach approximately 3.6594\% and 3.0800\%, respectively. Their combined improvement is not simply the sum of separate gains because the probability fields and overlap change jointly. The full union formulation accounts for these interactions. Redundancy is neither assumed absent nor presented as an algorithm-specific disadvantage; overlapping detection opportunities are intrinsic to the network objective.

\subsubsection*{Comparative Performance Analysis}
Figure~\ref{fig:algorithms} and Supplementary Table S1 compare measured runs under matched conditions. Powell reaches a mean fine-grid coverage near 4.34388\% over 20 initializations, with a mean of 327.1 evaluations from the shared allowance of 600. PSO reaches approximately 4.34360\% using 600 calls per run. Their final coverage values are very close; the practical distinction in this benchmark is the number of evaluations and the spread across starts, rather than a large coverage advantage over PSO.

The specified AFSA variant reaches approximately 2.5337\% on average, with substantial run-to-run variability, while uniform random search reaches approximately 2.9313\%. These outcomes expose the cost of local exploration and parameter choices in a bounded, small-budget comparison. They do not establish that AFSA in general is inferior to random search or to Powell. No algorithm-specific parameter tuning study was conducted, and the four-dimensional structured grid is substantially simpler than unrestricted deployment of many independently positioned sensors.

Median measured optimization times are approximately 1.57 s for Powell, 2.74 s for PSO, 2.82 s for the specified AFSA and 2.31 s for random search on the recorded system. Runtime is a computational metric and cannot be converted into deployment energy or real-time actuation capability. Supplementary Table S6 compares demonstrated evidence and remaining requirements for each deployment capability.

\begin{figure}[htbp]\centering
\includegraphics[width=\linewidth]{figures/fig6_algorithms.pdf}
\caption{\textbf{Repeated algorithm comparison under a common evaluation allowance.} (a) Median best training coverage and interquartile range across 20 runs per method. Early-stopped Powell traces are carried forward for visualization. (b) Distribution of final coverage evaluated on the fine rule. (c) Measured optimizer runtime, excluding final refined evaluation. Points represent individual runs; boxes show quartiles and whiskers use the standard 1.5-interquartile-range convention. AFSA denotes the explicitly specified variant.}\label{fig:algorithms}
\end{figure}
\FloatBarrier


\subsection*{Real-Ocean Current Replay and Environmental Transferability}
Figure~\ref{fig:ocean} connects the real current distributions to daily coverage in the held-out year. The fixed geometry, synthetic optimum and selected current-aware Powell geometry use the same nine nodes, sensing radius, target volume, residual law and physical constraints. Consequently their differences reflect layout choices under the declared environmental replay, rather than sensor-count growth or different detection models. The chronological split also avoids selecting a deployment from its reported test performance.

For South China Sea, held-out mean coverage is \OceanSCSFixed\% for the fixed grid, \OceanSCSSynthetic\% for the synthetic optimum and \OceanSCSOpt\% for current-aware Powell. The latter gains \OceanSCSGain\ percentage points over the fixed grid, with a temporal 95\% gain interval of [\OceanSCSGainLow,\OceanSCSGainHigh] percentage points. Corresponding Kuroshio-vicinity values are \OceanKuroFixed\%, \OceanKuroSynthetic\% and \OceanKuroOpt\%, and the gain interval is [\OceanKuroGainLow,\OceanKuroGainHigh] percentage points. For North Atlantic the coverages are \OceanNAFixed\%, \OceanNASynthetic\% and \OceanNAOpt\%, with a gain interval of [\OceanNAGainLow,\OceanNAGainHigh] percentage points. Supplementary Tables S7--S12 give exact inputs, held-out means, independent Monte Carlo intervals, repeated-method summaries, selected coordinates and state-reduction diagnostics.

The optimizer distributions and traces are retained for all 240 real-current runs (Supplementary Fig. S7), including poor starting points and variability. The selected-layout test comparison answers how a planning-year choice transfers to a later record; it does not establish that the selected method universally outperforms other optimizers. Monte Carlo and direct conditional integration agree within the archived sampling and numerical checks. State-reduction approximation, optimizer variability, integration accuracy and environmental variation remain separately reported.

The Kuroshio-vicinity result illustrates a transfer limitation in its point estimate: current-aware Powell reaches \OceanKuroOpt\% in the held-out year, compared with \OceanKuroSynthetic\% for the synthetic optimum. Improving the planning-year objective therefore does not guarantee improvement over the static optimum in a later year. Paired temporal gain intervals against the synthetic optimum include zero for South China Sea and Kuroshio, whereas the North Atlantic interval is positive (Supplementary Table S13). These intervals describe the empirical test record under the declared model; they do not establish general field superiority. The comparison retains the unfavorable Kuroshio point estimate and the uncertainty of the other comparisons.

Supplementary Fig. S6 varies the assumed horizon from 0 to 1, 5, 10 and 20 minutes and horizontal residual scales from 30 to 180 m, retaining the selected nominal layouts. These are stress tests, not test-year tuning. A zero horizon recovers the static model. Longer horizons increase the role of coherent drift and the limitations of a single fixed plan. Variations in daily forcing and residual assumptions can change the apparent advantage, so deployment guidance must name its horizon and uncertainty law.

The public derivative represents large-scale environmental forcing, approximately 100 km horizontally and daily in time. It does not resolve turbulence in a 100 m target volume, short-period tides or mooring response. It supplies a real ocean-data applicability check for the computational framework, while sensor-specific position errors, acoustic detection radius, local bathymetry, actuation and high-frequency flow require site calibration and field validation. The three sites are examples of distinct forcing records rather than a comprehensive validation of ocean conditions.

\begin{figure}[htbp]\centering
\includegraphics[width=\linewidth]{figures/fig7_real_ocean.pdf}
\caption{\textbf{Real Copernicus Marine currents and held-out coverage replay.} (a--c) Eastward and northward velocities at 150 m from the public $1^\circ$ GLORYS12V1 derivative; light points denote planning year 2018 and dark points held-out year 2019. (d--f) Daily conditional expected coverage in 2019 for the fixed grid, synthetic optimum and current-aware Powell layout selected from 2018 only. Coverage uses a 600 s frozen-current displacement assumption and parent residual scales $(120,120,40)$ m. Current velocities are real reanalysis data; displacement and coverage are model-derived quantities. Daily curves are unsmoothed.}\label{fig:ocean}
\end{figure}
\FloatBarrier

\section*{Discussion}
The corrected formulation changes the interpretation of cooperative coverage. Sensor independence permits a product of complements at one fixed target, but does not remove the spatial dependence among probability fields. Keeping the union inside the spatial integral is therefore essential. The resulting implementation provides three mutually checkable views of performance: an integral over a physical sensing sphere, an average over the target volume, and an independent Bernoulli simulation. Agreement among them is more informative than agreement between two calculations that share the same aggregation error.

The optimized layout is also inseparable from the physical depth model. Truncation at the shallow boundary shifts actual depth away from the nominal parent centre. This leads to a lower-bound solution in the current benchmark and explains why that solution is sensitive to a change in vertical uncertainty. Future experiments should estimate the error model from independent calibration trajectories, distinguish bias from dispersion, and evaluate deployment selected from one dataset on another. The present sensitivity analysis demonstrates the need for such calibration rather than asserting general robustness.

\textbf{CFD--LSTM coupling and environmental dynamics.} Environmental modeling offers an explicit extension: CFD prediction of underwater-vehicle dynamics\cite{liu2019}, data-driven fluid modeling\cite{brunton2020,vinuesa2023}, turbulence modeling\cite{novati2021}, glider-position forecasting\cite{qiu2023}, maritime trajectory prediction\cite{li2023,murray2022} and LSTM storm-surge correction\cite{giaremis2024} illustrate relevant but distinct tasks. A future coupling could use CFD to supply flow-dependent trajectories, use an LSTM to forecast residual displacement distributions, and pass calibrated time-dependent means and covariance parameters to the coverage evaluator. Such coupling requires specified flow boundary conditions, mesh and time-step convergence, training and held-out sequences, forecast horizons and probabilistic calibration. None of those components is inferred from the present Gaussian benchmark, and no CFD--LSTM performance score is claimed here.

\textbf{Dynamic deployment and a movement proxy.} To make the adaptive direction more concrete, Supplementary Fig. S2 adds a reproducible 13-step synthetic horizontal-bias scenario. A known bias $\mathbf b_k=(80\sin(k\pi/6),40\sin(k\pi/3))$ m is applied to the parent horizontal centres. Static nominal positions are compared with ideal compensation $\mathbf s_{i,k}=\mathbf s_i^*-(b_{kx},b_{ky},0)$. Perfect compensation preserves the actual parent centres by construction and consequently preserves coverage. This is a controlled limiting case, not a trained predictor, online estimator or closed-loop ocean trial. Its value is to expose the movement that compensation requires while quantifying the loss from ignoring known bias.

Total nominal travel $D=\sum_{k,i}\|\mathbf s_{i,k}-\mathbf s_{i,k-1}\|_2$ is recorded alongside coverage. It is a movement proxy in metres, not energy in joules. An energy-aware extension would require propulsion or buoyancy-actuation models, communication and sensing loads, standby consumption and an operating horizon, for example through $\max C-\lambda E$ or an energy-budget constraint. Work on wave-powered monitoring\cite{lu2024} concerns a physical energy system and cannot substantiate a numerical energy advantage for a placement optimizer. The comparative capability table therefore retains the energy topic while making its present evidential status explicit.

\textbf{Real-time operation and digital-twin integration.} Underwater robotic tracking\cite{masmitja2023} and digital-twin monitoring\cite{dong2025,mauro2023} motivate connections between estimation and control. However, a short offline optimizer runtime alone does not demonstrate real-time operation. Replanning would have to complete within a declared sensing-update period and account for actuation limits, communication latency, forecast error and correlated disturbances. The current optimizer is an offline planning component. The bias experiment provides a reproducible starting point for developing and testing that broader loop.

Several limitations qualify the findings. The spherical binary sensor model neglects acoustic attenuation, occlusion, directionality and variable detection probability. The physical box and error scales are assumed rather than calibrated. The synthetic benchmark assumes independent sensors; the real-current replay includes shared large-scale drift but still lacks calibrated small-scale residual correlations. Nominal separation does not ensure separation of realized positions. Only four grid parameters are optimized, and no global-optimality certificate is given. The benchmark is symmetric and strongly uncertain, favoring smooth probability fields. Population algorithms could perform differently in larger, irregular or multimodal deployment problems or with different budgets. The public $1^\circ$ daily ocean record cannot resolve metre-scale deployment dynamics. Finally, numerical and Monte Carlo agreement validates computation under the selected model, while field usefulness requires site-specific position, sensing and control measurements.

\section*{Conclusions}
This study develops a consistent three-dimensional probabilistic coverage framework for underwater sensing under localization uncertainty. Actual positions are constrained physically, nominal depth is distinguished from additive error, and network union probabilities are formed pointwise before spatial averaging. Spherical integration and independent Monte Carlo trials estimate the same expected covered-volume fraction.

In the declared nine-node benchmark, optimizing grid centre, pitch and common nominal depth increases coverage from \FixedCov\% to \OptCov\%. Controlled ablations separate pitch and depth effects, and repeated PSO runs reach nearly the same final coverage as Powell. Powell uses fewer evaluations under the specified allowance, while the tested AFSA variant shows greater variability. Sensitivity experiments reveal that the selected boundary layout can lose its advantage when the vertical error law changes. These findings support a reproducible static planning method and identify the calibration, dynamics and energy-modeling work required for practical adaptive deployment.

The real-current replay extends the static benchmark to observation-constrained ocean forcing with training-only planning and a held-out year. Conditional union probabilities preserve dependence induced by common currents. The resulting applicability evidence remains limited by spatial and temporal data resolution and by uncalibrated displacement and residual-position assumptions.

\section*{Data Availability}
Real environmental inputs are Copernicus Marine GLORYS12V1 reanalysis (product DOI 10.48670/moi-00021), accessed through Mercator Ocean International's public OceanBench $1^\circ$ historical archive\cite{cmems_glorys,oceanbench_archive}. The package includes three retained-depth NetCDF subsets, 2,190 extracted daily records, original metadata, exact requests and SHA-256 hashes. The derived replay outputs, all synthetic-benchmark results, raw bit-packed Monte Carlo outcomes, parameters, optimizer traces, maps, sensitivity tables and plotting scripts are included for offline reproduction. Source: European Union, Copernicus Marine Service Information or Data (2018--2019), \copyright Mercator Ocean. Original manuscript assets are retained locally as a revision record.

\section*{Code Availability}
The complete executable source, formula-to-code registry, pseudocode, frozen inputs and reproduction instructions are openly available at \url{https://github.com/houfence-design/underwater-coverage-reproducibility}. The source release distinguishes synthetic computations, real ocean inputs and model-derived replay results. \textbf{Before submission: publish the prepared release in a recognized DOI-assigning repository and insert the issued code-release DOI here.} The external Copernicus product DOI does not identify this project's code.

\section*{Acknowledgments}
The authors acknowledge the academic guidance and supervision of Kong Sibo. All revised scientific figures were generated from the accompanying computational scripts. This study has been conducted using E.U. Copernicus Marine Service Information: https://doi.org/10.48670/moi-00021. Mercator Ocean International provides the public historical derivative through OceanBench.
\section*{Author Contributions}
Kong Sibo and Hou Xinyu contributed to conceptualization and methodology. Hou Xinyu contributed to software, formal analysis, investigation, data curation, original drafting and project administration. Kong Sibo contributed resources, supervision and review and editing. \textbf{Author completion required: confirm the remaining validation and visualization roles for all five named authors and resolve the unmatched initials ``T.Z.'' in the original contribution statement before submission.}
\section*{Funding}
This research received no external funding.
\section*{Competing Interests}
The authors declare no competing interests.
\bibliography{references}
\end{document}
